\documentclass[11pt,a4paper]{article}
\usepackage[margin=0.78in,headheight=14pt,footskip=24pt]{geometry}
\usepackage[T1]{fontenc}
\usepackage{lmodern}
\usepackage{xcolor,listings,graphicx,float,fancyhdr,caption,titlesec,microtype}
\definecolor{codebg}{RGB}{247,247,247}
\definecolor{codegreen}{RGB}{0,110,55}
\definecolor{codeblue}{RGB}{30,70,160}
\definecolor{codepurple}{RGB}{135,40,140}
\lstdefinestyle{terminalstyle}{basicstyle=\ttfamily\scriptsize,backgroundcolor=\color{codebg},breaklines=true,frame=single,rulecolor=\color{black!18},columns=fullflexible,keepspaces=true,showstringspaces=false}
\lstdefinestyle{pythonstyle}{language=Python,basicstyle=\ttfamily\scriptsize,backgroundcolor=\color{codebg},keywordstyle=\color{codeblue}\bfseries,commentstyle=\color{codegreen}\itshape,stringstyle=\color{codepurple},breaklines=true,frame=single,rulecolor=\color{black!18},numbers=left,numberstyle=\tiny\color{gray},numbersep=6pt,columns=fullflexible,showstringspaces=false}
\pagestyle{fancy}
\fancyhf{}
\lhead{\small Big Data Analysis Lab}
\rfoot{\small Experiment 12}
\cfoot{\thepage}
\renewcommand{\headrulewidth}{0.4pt}
\setlength{\parindent}{0pt}
\setlength{\parskip}{4pt}
\titleformat{\section}{\large\bfseries}{\thesection}{0.7em}{}
\titleformat{\subsection}{\normalsize\bfseries}{\thesubsection}{0.7em}{}
\captionsetup{font=small,labelfont=bf}
\setlength{\intextsep}{5pt}
\setlength{\textfloatsep}{6pt}
\setlength{\abovecaptionskip}{3pt}
\setlength{\belowcaptionskip}{0pt}
\graphicspath{{images/}}
\begin{document}
\noindent
\textbf{Name:} Nikhil Sunchu\\
\textbf{Class:} D20A-31

\vspace{0.5em}
\begin{center}
{\Large\bfseries Experiment 12: PageRank}\\
\textbf{Implementation of PageRank using PySpark}
\end{center}

\section{Implementation Steps}

\subsection{Install PySpark}
Google Colab was used to implement the PageRank algorithm. PySpark was installed using pip. The installation was already available in the runtime and its requirement was verified.
\begin{lstlisting}[style=terminalstyle]
!pip install pyspark
\end{lstlisting}
\begin{figure}[H]
\centering
\includegraphics[width=0.98\linewidth]{install.png}
\caption{PySpark installation and package verification in Google Colab.}
\end{figure}

\subsection{Create the Input Dataset}
A text file named \texttt{pgdata.txt} was created with directed links between four web pages. Each line contains the source page followed by the destination page.
\begin{lstlisting}[style=terminalstyle]
page1 page3
page2 page1
page4 page1
page3 page1
page4 page2
page3 page4
\end{lstlisting}
\begin{figure}[H]
\centering
\includegraphics[width=0.66\linewidth]{dataset.png}
\caption{Input dataset containing the directed page links.}
\end{figure}

\subsection{Import Libraries and Initialize SparkContext}
The required PySpark classes were imported and a SparkContext was initialized with the application name \texttt{pgrank}.
\begin{lstlisting}[style=pythonstyle]
from pyspark import SparkContext, SQLContext

sc = SparkContext(appName="pgrank")

print("SparkContext initialized successfully!")
print("Application Name:", sc.appName)
\end{lstlisting}
\begin{figure}[H]
\centering
\includegraphics[width=0.98\linewidth]{spark_context.png}
\caption{SparkContext initialization and application name.}
\end{figure}

\subsection{Upload the Dataset}
The input file was uploaded to the Colab runtime using the file upload utility.
\begin{lstlisting}[style=pythonstyle]
from google.colab import files

linkfile = files.upload()
\end{lstlisting}
\begin{figure}[H]
\centering
\includegraphics[width=0.98\linewidth]{upload.png}
\caption{Uploading \texttt{pgdata.txt} to the Google Colab runtime.}
\end{figure}

\subsection{Read and Group the Page Links}
The text file was read using Spark. Each line was split into a source and destination page, duplicate links were removed, and destination pages were grouped by their source page.
\begin{lstlisting}[style=pythonstyle]
links = (
    sc.textFile("pgdata.txt")
    .map(lambda line: line.split())
    .map(lambda pages: (pages[0], pages[1]))
    .distinct()
    .groupByKey()
    .persist()
)

print("Page links loaded successfully:")
for page, neighbours in links.collect():
    print(page, list(neighbours))
\end{lstlisting}
\begin{figure}[H]
\centering
\includegraphics[width=0.98\linewidth]{links.png}
\caption{Grouped source pages and their outgoing links.}
\end{figure}

\subsection{Define the Contribution Function and Initialize Ranks}
The contribution function distributes a page's rank equally among its outgoing neighbours. Initially, every page was assigned a rank of 1.0.
\begin{lstlisting}[style=pythonstyle]
def computeContribs(neighbours, rank):
    for neighbour in neighbours:
        yield (neighbour, rank / len(neighbours))

ranks = links.map(lambda x: (x[0], 1.0))

print("Initial PageRank values:")
for pair in ranks.collect():
    print(pair)
\end{lstlisting}
\begin{figure}[H]
\centering
\includegraphics[width=0.98\linewidth]{initial_ranks.png}
\caption{Contribution function and initial PageRank values for the four pages.}
\end{figure}

\subsection{Execute the PageRank Iterations}
The PageRank values were updated for ten iterations. The rank contributions from incoming links were added using \texttt{reduceByKey()}, and the damping factor of 0.85 was applied in each iteration.
\begin{lstlisting}[style=pythonstyle]
n = 10

for x in range(n):
    contribs = links.join(ranks).flatMap(
        lambda pair: computeContribs(
            pair[1][0], pair[1][1]
        )
    )

    ranks = (
        contribs
        .reduceByKey(lambda v1, v2: v1 + v2)
        .map(lambda pair: (pair[0], pair[1] * 0.85 + 0.15))
    )

    print("Iteration", x)
    for pair in ranks.collect():
        print(pair)
\end{lstlisting}
\begin{figure}[H]
\centering
\includegraphics[width=0.90\linewidth]{iterations_code.png}
\caption{PageRank iteration code executed in Google Colab.}
\end{figure}
\begin{figure}[H]
\centering
\includegraphics[height=6.0in,keepaspectratio]{iterations_output.png}
\caption{PageRank values printed for iterations 0 through 9.}
\end{figure}

\end{document}
